\section{Gravity as Curvature}

Gravity gets its own chapter because it is special. The other forces are twists living on the crystal. Gravity is the bending of the crystal itself. This matches Einstein's great insight, that gravity is not a force pulling through space but the curving of space by what sits in it, and the crystal makes that idea literal.

\subsection{Matter bends the web}

The crystal is not perfectly uniform once matter is in it. A dense knot of pattern, a lump of matter, distorts the connections around it, changing the step-counting distances nearby. Other patterns moving through that distorted region follow the cheapest paths available, and those paths now bend toward the lump. From the outside it looks as if the lump is pulling them. There is no pull. There is only the bent web, and things taking the shortest road across it. That is gravity.

So gravity is the curvature of the crystal, responding to the matter in it, and matter responding back to the curvature. This is exactly Einstein's picture, here built from the web rather than assumed.

\subsection{It gives back the right laws}

The test of this is whether it reproduces known gravity, and the simulation walks up the ladder.

In the everyday, slow, weak case, it returns Newton's law: the familiar inverse-square pull, weaker with the square of the distance, that governs falling apples and planetary orbits. Then, beyond Newton, it yields the heart of Einstein's equation, a law relating the bending of the web to the matter and energy in it, that conserves energy properly, reduces to Newton when things are slow, and carries gravitational waves, ripples in the web, at the speed of light. The full nonlinear version holds in the form used for the whole expanding universe, with its parts staying consistent.

So from ``matter bends the web and things take cheap paths,'' the actual mathematics of gravity, Newton at low energy and Einstein beyond, comes out rather than being inserted.

\subsection{The graviton}

Light is a ripple in the twist-field, and the smallest such ripple is a photon. Gravity is a ripple in the web's curvature, and the smallest such ripple is a graviton. The crystal hosts this too. Built directly on the discrete web, the gravity ripple comes out with exactly the right credentials: it carries no mass, it travels at the speed limit, and it has the particular two-faced twist (spin two) that the carrier of gravity must have. The quantum of gravity, which ordinary physics struggles to make sense of, is just another pattern the crystal naturally supports.

\subsection{No infinities}

Ordinary gravity has a notorious disease: at the center of a black hole, or at the very beginning of the universe, the math says the curvature becomes infinite, a singularity, where physics breaks down. The crystal cures this for free. Because space is a discrete web with a smallest cell, there is a smallest possible length, and below it ``curvature'' stops meaning anything. So curvature can grow huge but never infinite. It is capped at a finite value set by the graininess. The singularity is softened into something finite and sensible. A world built from counting cannot have a true infinity in it, and that turns out to be exactly what gravity needed.

\textbf{Takeaway:} gravity is not a force but the bending of the crystal by matter, with things following the cheapest paths across the bent web, just as Einstein said. The crystal returns Newton's law, the core of Einstein's equation, gravitational waves, and a proper massless spin-two graviton, and its smallest length caps curvature so the dreaded singularities become finite.
